% Research-agenda addition (control strand): effective reviewers.
% Inserted right after Theorem~\ref{thm:assurance}.
\begin{remark}
\textbf{Effective reviewers.}
\pdstatus{\Vision}{}
The independence hypothesis of
Theorem~\ref{thm:assurance} is the first to fail when the reviewers are
language models. Kohli measured a panel of nine frontier models from seven
model families and found that the nine judges carried about two independent
votes' worth of information~\pdcite{kohli2026ninejudges}. The tasks were
classification and pairwise-preference judging, not flaw finding, so the
number transfers as a warning, not as a parameter. Family diversity alone did
not restore independence in that panel.

At this section's $d=0.6$, nine independent reviewers would miss a fixed flaw
with probability $0.4^9\approx2.6\times10^{-4}$. If the panel behaved like two
independent reviewers, it would miss the flaw with probability $0.4^2=0.16$,
about $610$ times as often. A panel can be summarized by an \emph{effective
number} measured on the panel itself: seed flaws with known answers, estimate
the all-miss rate $\hat m$ and the per-reviewer rate $\hat d$, and set
$k_{\mathrm{eff}}=\log\hat m/\log(1-\hat d)$. It uses the joint event the
bound needs; a pairwise statistic cannot replace it, by the counterexample in
\S\ref{sec:assurance-correlation}. But it describes the seeded flaws and
nothing more. Put back into the bound, it returns the measured average
all-miss rate, not a bound for any particular flaw, and a panel can miss one
kind of flaw every time while catching the rest. An assurance claim that uses
it must report its confidence interval and state the assumption that future
flaws are drawn like the seeded ones.
\end{remark}

\pdexercisepointerone{\ref{ex:he-effective-reviewers}}{\pageref{ex:he-effective-reviewers}}
